We are delighted to welcome you to MAGMaR 2026, the second workshop on Multimodal Augmented Generation via Multimodal Retrieval. MAGMaR is being held in San Diego, USA on the 4th of July, 2026, and is co-located with ACL 2026, which takes place from July 2nd through the 7th. This workshop is organized in support of ACL's Special Interest Group on Image and Language (SIGIL).

$ $\par

Audiovisual media is becoming an increasingly dominant form of online information consumption. From firsthand, ``in the wild'' video footage of natural disasters to professionally edited news coverage of major political events, videos serve as rich sources of information for producing factual, grounded articles. Especially for actively unfolding events, grounding articles in video can help combat misinformation and provide journalists and analysts with tools to quickly synthesize new developments.

$ $\par

Individual research groups have independently begun addressing this challenge, leading to parallel yet disconnected efforts to define the research space. ACL 2025 hosted the first MAGMaR workshop focused on Video Event Retrieval. This year's iteration focused on two primary areas: (1) the retrieval of multimodal content spanning text, images, audio, and video; and (2) retrieval-augmented generation, with an emphasis on multimodal retrieval and grounded generation. Relevant topics to the workshop this year included document retrieval, multimodal retrieval, retrieval-augmented generation (RAG), multimodal RAG, multimodal question answering, and research on video, image, and audio understanding.

$ $\par

To further this goal, we again hosted a shared task focused on video retrieval, and moreover, extended the task this year to include article generation from multiple videos. Specifically, it focused on retrieving relevant videos and generating grounded reports that respond to information needs. Given a query describing a real-world current event, participating systems needed to identify pertinent videos from a large multilingual, multimodal collection and use that evidence to produce a coherent and informative written report.

$ $\par

There were two tracks:

\begin{itemize}
    \item Retrieval: Systems provided a ranked list of videos in the collection ordered by relevance to the query.
    \item Generation: Systems produced a text report that answers the information need and grounds its content in the retrieved videos.
Teams were able to submit to either track or both.
\end{itemize}

$ $\par

We saw a large increase in the number of submissions to our shared task this year, with four teams submitting dozens of systems. All teams had at least one system that beat a very strong baseline and yielded some very interesting insights on what works and where are the open problems in this challenging multimodal domain. Check out the findings paper and teams' system descriptions for some really interesting analysis of how to build strong Video RAG systems.

%$ $\par
%
%Baseline models were OmniEmbed (with RankVideo reranking) for retrieval and CAG for generation. OmniEmbed was the winner of MAGMaR 2025's shared task which is why it was chosen for this year's baseline. The Retrieval track was evaluated using standard IR metrics: nDCG, Recall, and MAP at various cutoffs. Our official evaluation used the average of these scores. All RAG and Oracle Generation leaderboard scores were evaluated with CLUE (Zhang et al., 2026) and are MiRAGE (Martin et al., 2025) in the reference setting. Human evaluation was done by 3 human annotators. The annotators gave a scalar 1-5 score for each article and picked the overall best from the systems for each query.
%
%
%$ $\par
%
%Several teams submitted strong systems to the leaderboard. The best-performing submission,
%
%$ $\par
%
%For the retrieval task, the OmniEmbed Baseline had an average score across all metrics of 0.500. We had 3 teams submit numerous systems that beat this baseline. Overall, HAIVLab was able to make an impressive improvement to get to 0.832. The best performing system for MARS was 0.759, and Mixedbread got 0.719. Again, all of these beat the very strong baseline of last year's winning system by a substantial margin.
%
%$ $\par
%
%For the generation task, using the oracle ranking list, we use human scoring to rank systems using a Likert Scale of 1-5. The baseline CAG system received 3.088. We had 4 teams submit systems to this track, 2 of which beat the baseline. MARS' best system got a score of 3.833 slightly edging out Pengkhil's 3.825. Additionally, HAIVLab did have some systems receive best summary votes on specific queries. Once more, we received very strong systems that beat a tough baseline.

$ $\par

This year, the program of MAGMaR includes two keynote talks, one presentation session, and one poster session. With an increse in submissions from last year, we were able to accept 15 out of 26 papers, for an overall acceptance rate of 58\%. Of these, six were accepted as oral presentations. Once more, we allowed for non-archival submissions which has led to some interesting papers published in other venues that are being presented at the workshop. The members of our Program Committee and Organizing Comittee did an excellent job in reviewing the submitted papers, and we thank them for their essential role in selecting the accepted papers and helping produce a high quality program for the conference.



$ $\par


A workshop requires the hard work of numerous people, both behind the scenes and those that you will see more prominently. First off, we want to say thank you to our two keynote speakers, Nanyun (Violet) Peng from UCLA and Chenliang Xu from the University of Rochester, who have agreed to give talks about multimodal problems. Dr. Peng's talk ``Towards Self-Improving Multimodal Models'' tackles multimodal reasoning and generation problems, while Dr. Xu's talk ``Multi-level Alignment in Audio-Visual Scene Generation and Learning'' looks at aligning representations across modalities. Both of these cover challenging problems focused on in this workshop and explored in our shared task. We appreciate their insights.


$ $\par

Additonally, we would be remiss to not mention the people who helped organize (and participated) in our shared task on retrieving events in videos. Our online leaderboard received numerous submissions and grew substantially over last year.


$ $\par

Finally, we thank all contributors, reviewers, and attendees who helped make MAGMaR 2026 possible. We hope you enjoy a day full of engaging talks, thought-provoking posters, and stimulating discussion.

$ $\par

$ $\par

Kenton Murray and Reno Kriz, Editors
